\chapter{Introduction}
\label{ch:introduction}
\section{Problem and motivation}
Security screening of baggage involves recognizing suspicious objects under
occlusion, overlap and variable clutter. Embedded threats within electronic
devices make this problem more difficult because the host device and the
threat region are distinct concepts. A device box alone cannot establish a
threat, and a detector confidence is not an operational safety decision.
The IEDXray benchmark supplies a relevant research setting for studying these
distinctions in X-ray imagery~\citep{takele2026iedxray}.

RayGuard investigates a practical path from an authorized scan to localized
model findings, operator review and a traceable result record. The project
also considers a broader setting in which component evidence is distributed
across multiple packages. That setting requires externally justified grouping
and component labels; it cannot be obtained by arbitrarily combining unrelated
benchmark images. The second advisor meeting prioritized collection for that
distributed work while retaining a simple existing-model GUI showcase.
\projectrecord{docs/meeting_minutes/SDP_Meeting_02_Minutes_Fall_2026.md}

\section{Objectives and scope}
The original proposal calls for a 500-sample dataset with positive and negative
examples, image preprocessing, contemporary object detection, unseen-data
evaluation and a GUI supporting control and recording. The report evaluates
these objectives against actual evidence rather than treating a working screen
as completion of the full research problem.

The immediate SDP I implementation objectives are to preserve an explicit
task and category map, execute a verified detector, display original-coordinate
boxes and scores, distinguish failures from empty results, and retain provenance
through review and export. The experimental objectives are to reproduce a
supplied checkpoint's evaluation and to compare untuned and tuned models under
a controlled training-only validation design. These experiments remain distinct
from reproducing the authors' full training procedure.

\begin{table}[htbp]
\centering\small
\caption{Scope boundaries used throughout this report.}
\label{tab:scope}
\begin{tabular}{p{0.21\textwidth}p{0.31\textwidth}p{0.35\textwidth}}
\toprule Workstream & Current evidence & Remaining acceptance\\\midrule
Single-scan prototype & Bounded real inference and implemented operator workflow & Human rehearsal, eligible accuracy evaluation and full P1 decisions\\
Research software & Shared contracts, evaluator and initialization helper integrated; training draft available & Real trainer connection, reviewed data policy and qualified pilot\\
Distributed collection & Sanitized inventory and catalogue targets & QC, physical/case labels, accepted counts and collection completion\\
Distributed inference & Interface and research direction documented & Component evidence, case ground truth, fusion and evaluation\\\bottomrule
\end{tabular}
\end{table}

\section{Research background}
Baggage threat detection has been surveyed for both 2D and 3D screening
\citep{velayudhan2023survey}. Dark-field X-ray imaging combined with neural
networks has been reported to discriminate threat materials from benign
materials in proof-of-concept experiments \citep{astolfo2022darkfield}, and large
X-ray benchmarks such as SIXray address prohibited-item discovery in overlapping
images \citep{miao2019sixray}. These sources motivate the problem but do not
establish this project's performance.
The literature offers several detector architectures rather than one universal
solution. Faster R-CNN shares learned features with a region-proposal network,
providing a representative two-stage approach~\citep{ren2015faster}.
RetinaNet's focal loss addresses the imbalance encountered in dense detection,
which motivates attention to hard and rare cases without proving transfer to
this X-ray task~\citep{lin2017focal}. DETR formulates detection as set prediction
with bipartite matching~\citep{carion2020detr}; Grounding DINO incorporates grounded
pretraining for open-set detection~\citep{liu2023grounding}.
These papers motivate design alternatives, not claims that those models were
trained or evaluated by this team.

YOLOv10 proposes consistent dual assignments and an efficient architecture for
end-to-end detection~\citep{wang2024yolov10}. Its available task-matched checkpoint
and executed project backend make YOLOv10-M a practical reference for RayGuard.
This selection is based on implementation evidence and project constraints;
it does not imply that it is the highest-scoring model in the IEDXray paper.

The IEDXray release covers device, generic-explosive and specific-explosive
detection, alongside a complete annotation export. The inspected paper reports
17,360 images with a published 12,224/5,136 train/test split. Its protocol describes
separation of physical configurations and concealment setups, while the local
audit identifies four exact image-byte overlaps across the supplied splits.
The published claim and local finding must both be retained.
\projectrecord{docs/iedxray-paper-review.md}

FALCON studies component grounding and compositional reasoning within X-ray
scenes~\citep{michael2026falcon}. The project inspected its paper and release as
a research option; the release observations recorded in the repository do not
establish a running multi-package solution. Grounded component presence and
plausible language do not prove physical connectivity or dangerous assembly.
\projectrecord{docs/research.md}

\section{Literature and patent search record}
A scoping search was made on 22 November 2026 to locate related work and patents. It was limited to three web queries and is not a complete patent landscape; patent databases were not queried by classification or claim text. Table~\ref{tab:search} lists the searches, and the sources were read at the level of their abstracts and snippets.

\begin{table}[htbp]
\centering
\caption{Scoping literature and patent search, accessed 22 November 2026.}\label{tab:search}
\begin{tabular}{p{0.28\textwidth}p{0.38\textwidth}p{0.28\textwidth}}
\toprule
Query focus & Result found & Limit\\
\midrule
YOLO detectors for concealed IEDs in X-ray & Academic YOLO studies on IED replicas and small threat objects; one USPTO baggage-inspection patent that names YOLO as a detector & Snippets only; patent claims, assignee and status not verified\\
Multi-view or multi-scan fusion of detections & Multi-view detection with epipolar constraints; dual-view feature fusion; a dual-view dataset with auxiliary-view networks & Academic results only; no patent claiming cross-image detection fusion was found\\
Dark-field and multi-energy material discrimination & Reported in the supplied benchmark and survey citations & Cited through supplied papers; not independently re-read\\
Patent-focused queries (three further web searches, 22 Nov 2026) & Older related patents only: multi-sensor explosives detection with a neural network (EP0476070A1, US5153439A); X-ray and transmitted/scattered-ray baggage screening (US5600700A); dual-energy CT threat resolution (US20040258199A1); two-scanner dual-energy inspection (US9915752). No patent found for deep-learning detection of explosives hidden in modified electronics or for cross-image detection fusion & Web search snippets, not a patent-database query by classification or claim text; claims, status and assignees not verified. A Google Patents or USPTO full-text search with CPC G01V 5/20 is still required\\
Google Patents records of three hits (fetched 22 Nov 2026) & US8885938B2, Detecting concealed threats (Analogic; filed 2008, granted 2014): laptop-type object classification, neural network one option among SVM and PCA, mainly CT; KR102063859B1, airport AI baggage screening (Deepnoid; filed 2019): deep-learning X-ray classification of prohibited items; JP6678192B2, firearm detection with a neural network (Tsinghua University; filed 2018): not explosives & Status shown on the pages as active; claim text summarized by a fetch tool, not read verbatim, so claims are unverified. Not a complete database search: one of two Google Patents queries rendered; USPTO full-text search not run\\
\bottomrule
\end{tabular}
\end{table}

\section{Method and research contribution}
The engineering method begins with requirements and evidence inspection,
followed by dataset auditing, strict interfaces, bounded real inference and
application integration. Experimental preparation separates model execution
from evaluation. COCO conventions provide a common annotation and evaluation
language~\citep{lin2014coco}; the actual source category IDs remain task-specific.
For controlled improvement, Optuna supplies a framework for defining and
pruning search trials~\citep{akiba2019optuna}. Search is planned, not completed.
The selected optimization alternatives include SGD and AdamW; decoupled weight
decay motivates treating optimizer settings explicitly~\citep{loshchilov2019adamw}.

The current contribution is an auditable prototype and research-software path,
with documented findings and limitations. No new detector accuracy, tuning gain,
novel fusion method or operational safety claim is made in this draft.

\section{Repositories and individual contributions}
The shared research repository is
\url{https://github.com/abdalrahman-ismaik/SDP-I-RayGuard}.
The separate GUI repository is
\url{https://github.com/abdalrahman-ismaik/RayGuard-App}.
The application pins shared engine code through a Git submodule. GUI source
changes belong in the application checkout; its ignored predecessor is a recovery
copy. Appendix~\ref{app:contributions} records contribution evidence and outstanding
student verification. AI-assisted implementation and documentation must be disclosed
accurately; assistance does not establish a student's independent understanding.

\pending{Student-confirmed contributions, sanitized relevant prompt excerpts,
formal project number, final roster and institutional cover template. A scoped patent
search is not evidenced; this draft makes no patentability or clearance claim.}

\section{Report organization}
Chapter~\ref{ch:system} defines requirements and interfaces.
Chapter~\ref{ch:theory} explains the engineering model and evaluation quantities.
Chapter~\ref{ch:design} compares alternatives and documents constraints.
Chapter~\ref{ch:implementation} describes actual implementation and verification.
Chapter~\ref{ch:conclusions} evaluates progress and remaining work.
The appendices provide project-management evidence, requirement coverage,
risk preparation and contribution records.
